% Options for packages loaded elsewhere
% Options for packages loaded elsewhere
\PassOptionsToPackage{unicode}{hyperref}
\PassOptionsToPackage{hyphens}{url}
\PassOptionsToPackage{dvipsnames,svgnames,x11names}{xcolor}
%
\documentclass[
]{article}
\usepackage{xcolor}
\usepackage[top=1in,bottom=1in,left=1in,right=1in]{geometry}
\usepackage{amsmath,amssymb}
\setcounter{secnumdepth}{5}
\usepackage{iftex}
\ifPDFTeX
  \usepackage[T1]{fontenc}
  \usepackage[utf8]{inputenc}
  \usepackage{textcomp} % provide euro and other symbols
\else % if luatex or xetex
  \usepackage{unicode-math} % this also loads fontspec
  \defaultfontfeatures{Scale=MatchLowercase}
  \defaultfontfeatures[\rmfamily]{Ligatures=TeX,Scale=1}
\fi
\usepackage{lmodern}
\ifPDFTeX\else
  % xetex/luatex font selection
\fi
% Use upquote if available, for straight quotes in verbatim environments
\IfFileExists{upquote.sty}{\usepackage{upquote}}{}
\IfFileExists{microtype.sty}{% use microtype if available
  \usepackage[]{microtype}
  \UseMicrotypeSet[protrusion]{basicmath} % disable protrusion for tt fonts
}{}
\makeatletter
\@ifundefined{KOMAClassName}{% if non-KOMA class
  \IfFileExists{parskip.sty}{%
    \usepackage{parskip}
  }{% else
    \setlength{\parindent}{0pt}
    \setlength{\parskip}{6pt plus 2pt minus 1pt}}
}{% if KOMA class
  \KOMAoptions{parskip=half}}
\makeatother
% Make \paragraph and \subparagraph free-standing
\makeatletter
\ifx\paragraph\undefined\else
  \let\oldparagraph\paragraph
  \renewcommand{\paragraph}{
    \@ifstar
      \xxxParagraphStar
      \xxxParagraphNoStar
  }
  \newcommand{\xxxParagraphStar}[1]{\oldparagraph*{#1}\mbox{}}
  \newcommand{\xxxParagraphNoStar}[1]{\oldparagraph{#1}\mbox{}}
\fi
\ifx\subparagraph\undefined\else
  \let\oldsubparagraph\subparagraph
  \renewcommand{\subparagraph}{
    \@ifstar
      \xxxSubParagraphStar
      \xxxSubParagraphNoStar
  }
  \newcommand{\xxxSubParagraphStar}[1]{\oldsubparagraph*{#1}\mbox{}}
  \newcommand{\xxxSubParagraphNoStar}[1]{\oldsubparagraph{#1}\mbox{}}
\fi
\makeatother


\usepackage{longtable,booktabs,array}
\newcounter{none} % for unnumbered tables
\usepackage{calc} % for calculating minipage widths
% Correct order of tables after \paragraph or \subparagraph
\usepackage{etoolbox}
\makeatletter
\patchcmd\longtable{\par}{\if@noskipsec\mbox{}\fi\par}{}{}
\makeatother
% Allow footnotes in longtable head/foot
\IfFileExists{footnotehyper.sty}{\usepackage{footnotehyper}}{\usepackage{footnote}}
\makesavenoteenv{longtable}
\usepackage{graphicx}
\makeatletter
\newsavebox\pandoc@box
\newcommand*\pandocbounded[1]{% scales image to fit in text height/width
  \sbox\pandoc@box{#1}%
  \Gscale@div\@tempa{\textheight}{\dimexpr\ht\pandoc@box+\dp\pandoc@box\relax}%
  \Gscale@div\@tempb{\linewidth}{\wd\pandoc@box}%
  \ifdim\@tempb\p@<\@tempa\p@\let\@tempa\@tempb\fi% select the smaller of both
  \ifdim\@tempa\p@<\p@\scalebox{\@tempa}{\usebox\pandoc@box}%
  \else\usebox{\pandoc@box}%
  \fi%
}
% Set default figure placement to htbp
\def\fps@figure{htbp}
\makeatother





\setlength{\emergencystretch}{3em} % prevent overfull lines

\providecommand{\tightlist}{%
  \setlength{\itemsep}{0pt}\setlength{\parskip}{0pt}}



 


\makeatletter
\@ifpackageloaded{caption}{}{\usepackage{caption}}
\AtBeginDocument{%
\ifdefined\contentsname
  \renewcommand*\contentsname{Table of contents}
\else
  \newcommand\contentsname{Table of contents}
\fi
\ifdefined\listfigurename
  \renewcommand*\listfigurename{List of Figures}
\else
  \newcommand\listfigurename{List of Figures}
\fi
\ifdefined\listtablename
  \renewcommand*\listtablename{List of Tables}
\else
  \newcommand\listtablename{List of Tables}
\fi
\ifdefined\figurename
  \renewcommand*\figurename{Figure}
\else
  \newcommand\figurename{Figure}
\fi
\ifdefined\tablename
  \renewcommand*\tablename{Table}
\else
  \newcommand\tablename{Table}
\fi
}
\@ifpackageloaded{float}{}{\usepackage{float}}
\floatstyle{ruled}
\@ifundefined{c@chapter}{\newfloat{codelisting}{h}{lop}}{\newfloat{codelisting}{h}{lop}[chapter]}
\floatname{codelisting}{Listing}
\newcommand*\listoflistings{\listof{codelisting}{List of Listings}}
\makeatother
\makeatletter
\makeatother
\makeatletter
\@ifpackageloaded{caption}{}{\usepackage{caption}}
\@ifpackageloaded{subcaption}{}{\usepackage{subcaption}}
\makeatother
\usepackage{bookmark}
\IfFileExists{xurl.sty}{\usepackage{xurl}}{} % add URL line breaks if available
\urlstyle{same}
% fallback for those not using the hyperref driver hyperxmp:
\makeatletter
\@ifundefined{xmpquote}{\newcommand{\xmpquote}[1]{#1}}{}
\makeatother
\hypersetup{
  pdftitle={Columbia River Spring Chinook Marine Survival: Synthesis of Multi-Stage SEM Deliverables},
  pdfauthor={Lisa Crozier et al.},
  colorlinks=true,
  linkcolor={blue},
  filecolor={Maroon},
  citecolor={Blue},
  urlcolor={Blue},
  pdfcreator={LaTeX via pandoc}}


\title{Columbia River Spring Chinook Marine Survival: Synthesis of
Multi-Stage SEM Deliverables}
\usepackage{etoolbox}
\makeatletter
\providecommand{\subtitle}[1]{% add subtitle to \maketitle
  \apptocmd{\@title}{\par {\large #1 \par}}{}{}
}
\makeatother
\subtitle{Consolidated Results Across Objectives 1--5}
\author{Lisa Crozier et al.}
\date{2026-11-06}
\begin{document}
\maketitle

\renewcommand*\contentsname{Table of contents}
{
\setcounter{tocdepth}{2}
\tableofcontents
}

\subsection{\#tinytex::tlmgr\_install(``koma-script'')}\label{tinytextlmgr_installkoma-script}

title: ``Columbia River Spring Chinook Marine Survival: Synthesis of
Multi-Stage SEM Deliverables'' subtitle: ``Consolidated Results Across
Objectives 1--5'' author: ``Lisa Crozier et al.'' date: today format:
pdf: pdf-engine: pdflatex documentclass: article \# latex-auto-install:
false toc: true toc-depth: 2 number-sections: true colorlinks: true
geometry: - top=1in - bottom=1in - left=1in - right=1in execute: echo:
false warning: false message: false ---

\newpage

\section{Objective 1 \& 2: Ecosystem Drivers \& Multi-Model Variable
Importance}\label{objective-1-2-ecosystem-drivers-multi-model-variable-importance}

\subsection{Overview}\label{overview}

Model selection across candidate single-indicator structural equation
models (SEMs) evaluated the relative influence of coastal Northern
California Current (NCC) and Gulf of Alaska (AK) drivers on early marine
survival (SAR). Multi-model inference via Akaike weights (\(w_i\)) was
calculated across all statistically valid models
(\(p(\chi^2) \ge 0.05\), \(\text{CFI} \ge 0.80\),
\(\text{SRMR} \le 0.10\), \(\text{AGFI} \ge 0.80\)) following baseline
exclusions (\texttt{sablefishAK\_biomass4plus} and
\texttt{pollockGOA\_SSB\_lag1}).

\subsection{Figure 1: Relative Variable Importance across Ecosystem
Indicators}\label{figure-1-relative-variable-importance-across-ecosystem-indicators}

Relative variable importance (\(w_i\)) across all ecosystem indicators
meeting the \(\max(w_i) \ge 0.05\) threshold, comparing sequential
mediation (DAG 1A) and direct growth mediation (DAG 1B) topologies
across the 24-year complete dataset and 19-year expanded series.

\begin{figure}

\centering{

\includegraphics[width=1\linewidth,height=\textheight,keepaspectratio]{Outputs/figures/obj2_fig1_variable_importance_barchart.pdf}

}

\caption{\label{fig-obj2-importance}Relative variable importance (Akaike
weights) across candidate functional ecosystem indicators.}

\end{figure}%

\newpage

\section{Objective 3: Direct vs.~Modulated Predation
Topologies}\label{objective-3-direct-vs.-modulated-predation-topologies}

\subsection{Overview}\label{overview-1}

To address counter-intuitive positive relationships between raw predator
abundance and salmon survival, candidate predators were evaluated under
an effective predation framework (\(P_{\text{eff}}\)):
\[P_{\text{eff}} = P_{\text{norm}} \cdot \text{SST}_{\text{gate}} \cdot (1 - \text{AltPrey}_{\text{gate}})\]
where thermal spatial overlap (\(\text{SST}\)) and alternative forage
buffering (\(\text{AltPrey}\)) were scaled with a 0.2 baseline floor to
prevent complete zero-multiplication.

\subsection{\texorpdfstring{Table 3: Effective Predation
(\(P_{\text{eff}}\)) Model
Comparisons}{Table 3: Effective Predation (P\_\{\textbackslash text\{eff\}\}) Model Comparisons}}\label{table-3-effective-predation-p_texteff-model-comparisons}

Comparison of nominal predator biomass models versus top-ranked
modulated \(P_{\text{eff}}\) formulations. Shaded green cells indicate
models where thermal gating and/or alternative prey buffering restored
ecologically coherent top-down control (\(\beta < 0\), \(p \le 0.05\))
with substantial AIC improvement (\(\Delta\text{AIC} \le -2.0\)).

\begin{table}

\caption{\label{tbl-obj3-peff}}

\centering{

\begin{figure}
\centering
\includegraphics[width=1\linewidth,height=\textheight,keepaspectratio,alt={Modulation of marine predator exposure: Nominal biomass vs.~effective predation exposure (Peff).}]{Outputs/tables/obj3_table1_peff_model_comparisons.png}
\caption{Modulation of marine predator exposure: Nominal biomass
vs.~effective predation exposure (Peff).}
\end{figure}

}

\end{table}%

\newpage

\section{Objective 4: Multi-Basin Life-Stage
Bottlenecks}\label{objective-4-multi-basin-life-stage-bottlenecks}

\subsection{Overview}\label{overview-2}

The life-stage hypothesis ladder was tested sequentially across the
24-year complete dataset (1998--2021) to evaluate the spatial and
temporal bounds of survival bottlenecks: 1. \textbf{Single-Arena (Early
Marine NCC Only):}
\(\text{SAR} \sim \text{Growth}_{\text{NCC}} + \text{CPUE}_{\text{NCC}}\)
2. \textbf{Dual-Arena (+ Subarctic Alaska Marine):}
\(\text{SAR} \sim \text{Growth}_{\text{NCC}} + \text{CPUE}_{\text{NCC}} + \text{Prey}_{\text{AK}} + P_{\text{ocean\_eff}}\)
3. \textbf{Triple-Arena (+ Adult Columbia River Pinnipeds):}
\(\text{SAR} \sim \text{Growth}_{\text{NCC}} + \text{CPUE}_{\text{NCC}} + \text{Prey}_{\text{AK}} + P_{\text{ocean\_eff}} + P_{\text{river\_adult}}\)

\#\# Table 4: Single-, Dual-, and Triple-Arena Model Comparisons

\begin{table}

\caption{\label{tbl-obj4-arenas}}

\centering{

\begin{figure}
\centering
\includegraphics[width=0.95\linewidth,height=\textheight,keepaspectratio,alt={Life-stage arena bottlenecks: Single- (coastal baseline), Dual- (adding subarctic marine), and Triple-Arena (adding adult in-river pinnipeds) SEM comparisons.}]{Outputs/tables/obj4_table1_ocean_plus_CR_predation_24yr.png}
\caption{Life-stage arena bottlenecks: Single- (coastal baseline), Dual-
(adding subarctic marine), and Triple-Arena (adding adult in-river
pinnipeds) SEM comparisons.}
\end{figure}

}

\end{table}%

\newpage

\subsection{Figure 2: Final Pruned Latent Structural Equation
Models}\label{figure-2-final-pruned-latent-structural-equation-models}

\subsubsection{Candidate 2: Subarctic Groundfish Piscivore (Sablefish
Age
3+)}\label{candidate-2-subarctic-groundfish-piscivore-sablefish-age-3}

Final pruned structural model retaining both coastal early marine and
subarctic marine stages.

\begin{figure}

\centering{

\includegraphics[width=0.9\linewidth,height=\textheight,keepaspectratio]{Outputs/figures/obj4_fig1A_sablefish_model.png}

}

\caption{\label{fig-obj4-sablefish}Final pruned structural equation
model: Sablefish Age 3+ Peff model (AIC = -23.22, SAR R² = 0.73).}

\end{figure}%

\newpage

\subsubsection{Candidate 4: Subarctic Apex Shark (Sleeper
Shark)}\label{candidate-4-subarctic-apex-shark-sleeper-shark}

Final pruned structural model retaining top-down subarctic apex shark
predation.

\begin{figure}

\centering{

\includegraphics[width=0.9\linewidth,height=\textheight,keepaspectratio]{Outputs/figures/obj4_fig1B_sleepershark_model.png}

}

\caption{\label{fig-obj4-sleepershark}Final pruned structural equation
model: Sleeper Shark Peff model.}

\end{figure}%

\newpage

\subsubsection{Candidate 1 / Baseline: Coastal-Only Pruned
Architecture}\label{candidate-1-baseline-coastal-only-pruned-architecture}

Pruned architecture resulting when candidate subarctic marine predators
drop out (\(p \ge 0.05\)).

\begin{figure}

\centering{

\includegraphics[width=0.9\linewidth,height=\textheight,keepaspectratio]{Outputs/figures/obj4_fig1C_coastal_baseline_model.png}

}

\caption{\label{fig-obj4-baseline}Final pruned baseline structural
architecture (Coastal-only survival; AIC = -21.58, SAR R² = 0.83).}

\end{figure}%

\newpage

\section{Supplementary Material}\label{supplementary-material}

\subsection{Table S1: Complete Variable Importance Across All Evaluated
Indicators}\label{table-s1-complete-variable-importance-across-all-evaluated-indicators}

Full numerical compilation of cumulative Akaike weights (\(w_i\)) across
all candidate indicators, including those falling below the main text
\(\max(w_i) \ge 0.05\) plotting threshold.

\begin{table}

\caption{\label{tbl-supp-varimp}}

\centering{

\begin{figure}
\centering
\includegraphics[width=0.9\linewidth,height=\textheight,keepaspectratio,alt={Supplemental Table S1: Full compilation of relative variable importance across all candidate ecosystem indicators.}]{Outputs/tables/obj2_tableS3_variable_importance_wide.png}
\caption{Supplemental Table S1: Full compilation of relative variable
importance across all candidate ecosystem indicators.}
\end{figure}

}

\end{table}%




\end{document}
